\section{Bring Home The Bacon}
Bring Home The Bacon was a byproduct prototype of the Master thesis written by
Henriksen in 2012\cite{BHTB}. In the thesis, Henriksen sought to prove that having
a central, graphical-oriented, system to administer multiple Snort sensors and
their rules results in a signficant improvement in user workflow over the more
traditional command line environment. 

Henriksens tests showed that such a system could double the users
work efficiency, if not more, as the command line environment would scale very
badly once the number of sensors increases. Henriksen also outlines a rule
management process, which we have taken into account in our initial work to map
the program workflow in section \ref{sec:program_workflow}.

Despite being a prototype, BHTB does offer the essential functionality required
for managing rulesets on multiple sensors. It can display a list of rules
where you can turn those on or off on different sensors. It can also display a
list of rulesets and you can turn those on or off, but that only propagates down
to the rules, turning those on or off. That has been a drawback for BHTB, as you
loose whatever preferred setup you had if you simply want to turn one ruleset on
or off for a short while, and thus have to remember how it was set up before.

BHTB also lets you add thresholds and suppressions, but the system does not do
any validation of the input, thus leaving the Snort sensor vulnerable to
crashes\footnote{Snort will simply not restart if something in the configuration
files is malformed.} if the user inputs something wrong.

The way BHTB handles distribution is to generate all the files centrally and
then send them to the sensors one by one via SCP. It determines which rules to
distribute by comparing the main rule table with a second table representing
which rules the sensor should not have. 

On the technical side, BHTB is programmed with Python and uses the framework
Tornado\cite{tornado} as its webserver, which works
similar to Django with it's URL API system which you can read about
in section \ref{kapittel om django urls}.
It also uses SQLite\cite{sqlite} as it's database, which has limited abilities
when it comes to relations and more complex data structures. 

BHTB really begins to show it's
limitations once the system begins to be populated by a lot of
rules\footnote{In the five digit range.} and many sensors, thus having a
scalability problem.
Avdeling BKI has used BHTB in a production environment since 2012 and experiences the software as slow and
cumbersome. 


